\documentclass[final,3p,times,twocolumn]{elsarticle}

\usepackage[utf8]{inputenc}
\usepackage{xspace}
\newcommand{\spin}{SPIN\xspace}
\usepackage{amsmath,amssymb,amsthm}
\usepackage{url}
\usepackage{booktabs}
\usepackage{enumitem}
\setlist[itemize]{itemsep=3pt, topsep=3pt, parsep=0pt, partopsep=0pt, leftmargin=*}

\newtheorem{proposition}{Proposition}
\newtheorem{remark}{Remark}

\journal{Journal of Process Control}

\begin{document}

\begin{frontmatter}

\title{Reading the Sensitivity Peak of a PI Loop\\
from a Single Step Response}

\begin{abstract}
The sensitivity peak $M_s$ tells an engineer how robust a loop is.
Measuring it normally needs a plant model or a frequency sweep, and most
installed loops have neither. This paper shows that a single ordinary
setpoint step already contains the answer, and gives a simple law for
reading it.

\begin{itemize}
\item \emph{Why is $M_s$ in the record?} For a unit setpoint step the
error transform is $E(s)=S(s)/s$. The difference of the record is
therefore the impulse response of the sensitivity function itself. No
second experiment is needed.

\item \emph{How is it read?} Count how many times the error curls around
its settling point. Call the count $N$. We derive the law
$M_s=N/(c\sqrt{n_m^2+n^2})$, where $c=\ln(1/\varepsilon)/2\pi$ is fixed
by the counting rule and $n,n_m$ are the loop's phase and magnitude
slopes where the phase passes $-180^\circ$. The distance to instability
cancels out of the derivation, which is what makes the reading possible
without a model.

\item \emph{Why does dead time not disturb it?} Because a dead time
contributes the same amount to phase and to phase slope. When the
integral time cancels the dominant lag, the crossover condition pins the
dead time's phase share, and therefore its slope share, at exactly
$\pi/2$---whatever the dead time is. We compute $n=\pi/2$ for relative
dead times from $0.1$ to $0.9$, to three digits.

\item \emph{Why does it fail for PID?} Derivative action supplies phase
lead, so the dead time no longer has to supply the whole $-\pi/2$. Its
share, and hence $n$, then depends on how much dead time there is. The
invariance is lost by construction, not by accident.
\end{itemize}

\noindent
Measurement confirms each prediction with no fitted constant. The
dimensionless ratio the law predicts to be unity is measured at $1.256$,
$1.131$ and $1.074$ at $M_s\approx3.5$, $6.5$ and $8.5$---converging as
the theory requires, with interquartile ranges a percent or two wide,
and agreeing across plant orders one to three. The map moves by at most
$0.064$ across relative dead times $0.3$ to $0.9$ for PI, against
$2.58$ for PID at $\tau=0.3$; and readings on fixed classical tunings
have median error $+5\%$, with $22$ of $24$ within $\pm25\%$. A validity
condition read from the same record---does the error change sign---
rejects every reading that would have been wrong.
\end{abstract}

\end{frontmatter}

\section{Introduction}

The sensitivity peak $M_s=\max_\omega|S(j\omega)|$, with $S=1/(1+CG)$,
is the standard measure of how robust a control loop is. A value near
$1.4$ is comfortable, a value near $2.0$ is marginal, and the design
window $M_s\in[1.2,2.0]$ is long established
\cite{AstromHagglund2004,Skogestad2003}.

Getting the number is the problem. It is defined in the frequency
domain, so it needs the loop transfer function: either a model, which
someone must identify, or a frequency sweep, which someone must
commission. Of the eleven thousand loops surveyed by Desborough and
Miller \cite{Desborough2002}, almost none have either, and the survey's
main finding is that many of them run badly.

What those loops do produce is setpoint steps. This paper shows that one
such step is enough. The argument is short:

\begin{enumerate}[label=(\roman*)]
\item The step record \emph{is} the sensitivity function, because
$E(s)=S(s)/s$.
\item Near the point where a loop becomes marginal, one pole pair
dominates. Its damping sets the height of the peak.
\item Counting how many times the error curls around its settling point
measures that damping.
\end{enumerate}

Putting these together gives a linear law relating the count to $M_s$
(Section~\ref{sec:law}). The law has one free constant, the loop's phase
slope at crossover. Section~\ref{sec:delay} shows why that constant is
fixed at $\pi/2$ for a PI loop no matter how much dead time the plant
has, which is the reason the method needs no plant knowledge.
Section~\ref{sec:pid} shows why the same argument fails once derivative
action is added. Sections~\ref{sec:valid} and~\ref{sec:measure} give the
validity condition and the measurements.

We restrict attention to PI control. This is where the result is clean,
and it is also where the loops are: derivative action is switched off in
most industrial installations \cite{AstromHagglund2001}.

The counting device is the turn index of \cite{SPIN}, used there as a
damping diagnostic with limits fixed by calibration. Giving it a
frequency-domain meaning turns it into a measurement and, as a
by-product, replaces that calibration by a specification
(Section~\ref{sec:spin}).

\section{The law}
\label{sec:law}

\subsection{The record is the sensitivity function}

For a unit setpoint step,
\begin{equation}
E(s)=\frac{S(s)}{s}.
\label{eq:ES}
\end{equation}
The recorded error carries all of $S$, and its first difference is the
impulse response of $S$. This is why one ordinary step suffices and no
dedicated experiment is required.

\subsection{Counting the curls}

Plot the error against its own difference and watch the trajectory spiral
into the point where the loop settles. Let $N$ be the number of turns it
makes, counted after normalizing each axis and stopping when the curve
enters a small disc of radius $\varepsilon$ about that point
\cite[Def.~1]{SPIN}.

A loop that is well damped barely completes a quarter turn. A loop close
to instability winds many times. For a mode $e^{-\alpha t}\cos\omega t$
the spiral is logarithmic, and the count is
\begin{equation}
N=c\,\frac{\omega}{\alpha},\qquad c=\frac{\ln(1/\varepsilon)}{2\pi},
\label{eq:count}
\end{equation}
which is \cite[Prop.~3]{SPIN}. With the default $\varepsilon=0.1$,
$c=0.3665$. The count is a pure number: it does not change if the step
is made larger or the clock runs faster.

\subsection{From count to peak}

Near the stability boundary the loop loses stability through a single
complex pole pair \cite[Prop.~2]{SPIN}, at the frequency
$\omega_{180}$ where the phase passes $-\pi$. Let
$m=1-|L(j\omega_{180})|$ say how close the loop is to that boundary,
and let
\begin{equation}
n=-\frac{d\arg L}{d\ln\omega},\qquad
n_m=-\frac{d\ln|L|}{d\ln\omega}
\label{eq:slopes}
\end{equation}
be the phase and magnitude slopes there, both dimensionless. Writing
$u=\ln(\omega/\omega_{180})$ and expanding in logarithmic coordinates,
\begin{equation}
L(j\omega)\approx-(1-m)\,e^{-n_mu}e^{-jnu},
\label{eq:expand}
\end{equation}
so that to first order $1+L\approx m+n_mu+jnu$. Minimising over $u$
gives the distance from the locus to the critical point, and hence
\begin{equation}
M_s=\frac{\sqrt{n_m^2+n^2}}{m\,n}.
\label{eq:Ms}
\end{equation}

The same expansion locates the pole. Solving $1+L=0$ near
$j\omega_{180}$ gives a pair with frequency $\omega\approx\omega_{180}$
and decay rate $\alpha=m\,\omega_{180}n/(n^2+n_m^2)$, so
\begin{equation}
\frac{\omega}{\alpha}=\frac{n^2+n_m^2}{m\,n}.
\label{eq:woa}
\end{equation}

Dividing \eqref{eq:woa} by \eqref{eq:Ms} eliminates $m$---the distance
to the boundary, which the record cannot report---and leaves, with
\eqref{eq:count},
\begin{equation}
\boxed{\;M_s=\frac{\omega/\alpha}{\sqrt{n_m^2+n^2}}
=\frac{N}{c\sqrt{n_m^2+n^2}}\;}
\label{eq:law}
\end{equation}

Three things are worth stating plainly. The relation is \emph{linear} in
the count. The only plant-dependent quantity in it is the slope pair at
one frequency: plant order, lag structure and dead time enter nowhere
else. And $m$ has cancelled, so no knowledge of the gain margin is
needed---which is what makes the reading possible from a record alone.

Equation~\eqref{eq:law} is a first-order result in $m$, hence exact in
the limit of a marginal loop. For a comfortable loop $|S|$ does not fall
below one far from resonance, and the baseline adds a correction that
decays as $M_s$ grows. Section~\ref{sec:measure} measures how fast.

\section{Why dead time does not matter}
\label{sec:delay}

Equation~\eqref{eq:law} would be of little use if $n$ varied from plant
to plant, since $n$ is not measurable from the step record. For PI
control it does not vary, and the reason is elementary.

A dead time $L_d$ contributes phase $-L_d\omega$ and phase slope
$-d(-L_d\omega)/d\ln\omega=L_d\omega$. \emph{Its slope contribution
equals its phase contribution.} A first-order lag does not share this
property: it contributes phase $-\arctan(T\omega)$, approaching $\pi/2$,
while its slope $T\omega/(1+T^2\omega^2)$ falls back to zero.

Now take a PI controller whose integral time cancels the dominant lag,
$T_i=T$, which is the IMC and SIMC choice \cite{Skogestad2003}. The loop
phase is then
\begin{equation}
\arg L(j\omega)=-\frac{\pi}{2}-L_d\omega ,
\label{eq:argL}
\end{equation}
the $-\pi/2$ coming from the integrator. The crossover condition
$\arg L=-\pi$ therefore forces
\begin{equation}
L_d\omega_{180}=\frac{\pi}{2},
\label{eq:pin}
\end{equation}
and since the dead time's slope equals its phase,
\begin{equation}
n=L_d\omega_{180}=\frac{\pi}{2}\quad\text{exactly, for every }L_d .
\label{eq:npi2}
\end{equation}

A longer dead time simply moves the crossover to a lower frequency. The
product $L_d\omega_{180}$, which is what $n$ depends on, does not move.
This is the whole reason the reading needs no plant knowledge.

Table~\ref{tab:nphi} confirms \eqref{eq:npi2} numerically, and gives
$n_m$, which sits at unity because the cancelled lag leaves the
integrator's single $-1$ slope.

\begin{table}[t]
\centering
\caption{Phase and magnitude slope at crossover, PI with $T_i=T$, and
the slope of the law \eqref{eq:law} they predict. $\tau=L_d/(L_d+T)$.}
\label{tab:nphi}
\begin{tabular}{lccccc}
\toprule
$\tau$ & 0.1 & 0.3 & 0.5 & 0.7 & 0.9 \\
\midrule
$n$   & 1.57 & 1.57 & 1.57 & 1.57 & 1.57 \\
$n_m$ & 1.00 & 1.00 & 1.00 & 1.00 & 1.00 \\
$\omega_{180}$ & 14.14 & 3.67 & 1.57 & 0.67 & 0.18 \\
predicted slope & 1.47 & 1.47 & 1.47 & 1.47 & 1.47 \\
\bottomrule
\end{tabular}
\end{table}

\begin{remark}[Scope]
\label{rem:scope}
Equation~\eqref{eq:npi2} rests on the cancellation $T_i=T$. Away from
it the arctangents no longer cancel and $n$ drifts: over
$T_i/T\in[0.5,8]$ and $\tau\in[0.1,0.9]$ it ranges from $1.14$ to
$2.62$, so the slope of \eqref{eq:law} varies by a factor of about two
in the worst case. The law itself is unaffected---only its constant is.
Section~\ref{sec:measure} measures what this costs in practice, which is
less than the range suggests, because tunings in use cluster near
cancellation.
\end{remark}

\section{Why derivative action breaks it}
\label{sec:pid}

The argument of Section~\ref{sec:delay} used one fact: the integrator
supplies $-\pi/2$ of phase, the lag is cancelled, and so the dead time
must supply the remaining $-\pi/2$. Its slope share is then pinned
because slope and phase coincide for a dead time.

Derivative action breaks the first step. It supplies phase \emph{lead},
so the dead time no longer has to supply $-\pi/2$; it supplies less, and
how much less depends on where the crossover falls, which depends on how
much dead time there is. The pinning argument therefore fails, and $n$
becomes a function of $\tau$:

\begin{center}
\begin{tabular}{lccccc}
\toprule
$\tau$ & 0.1 & 0.3 & 0.5 & 0.7 & 0.9 \\
\midrule
$n$, PI ($T_d=0$)   & 1.57 & 1.57 & 1.57 & 1.57 & 1.57 \\
$n$, PID ($T_d/T_i=0.25$) & 2.82 & 2.28 & 1.47 & 1.48 & 1.57 \\
\bottomrule
\end{tabular}
\end{center}

Across $\tau\in[0.3,0.9]$ the PID phase slope spreads by $48\%$, against
$0\%$ for PI. The failure is not a numerical accident: it follows
directly from removing the condition that pinned $n$.

The theory also says where PID should still work. At large $\tau$ the
plant's dead time dominates the derivative's lead, and the spread falls
to $6\%$ over $\tau\in[0.5,0.9]$. Measurement agrees: recalibrating the
map with $T_d/T_i=0.25$ shifts it by $0.254$ at $\tau=0.9$ and $0.314$
at $\tau=0.7$---tolerable---but by $2.583$ at $\tau=0.3$. A PID reading
is therefore usable on delay-dominant loops and not on balanced ones,
exactly as \eqref{eq:npi2} predicts.

Strong derivative action fails for a second reason: the controller zeros
capture the critical pole pair, so the single-dominant-pair picture
behind \eqref{eq:expand} no longer holds. This is the mechanism already
noted in \cite[Rem.~3]{SPIN}, and it is why the map degrades sharply
past $T_d/T_i=0.25$.

\section{When the reading is valid}
\label{sec:valid}

The expansion \eqref{eq:expand} assumes the dominant pair is visible in
the record. Sometimes it is not. If integral action is very slow, the record
is dominated by a slow, non-oscillating decay, and the resonant mode
rides on it as ripple that never crosses zero. The trajectory then stays
on one side of the settling point and cannot wind, whatever the true
$M_s$ is.

Because the count is essentially half the number of times the error
crosses zero, the condition for a valid reading is as simple as it could
be:

\begin{quote}
\emph{The reading is valid when the error record changes sign.}
\end{quote}

No threshold, no parameter, no model. When the record is one-signed the
instrument returns nothing rather than something wrong.

Two facts make this workable. Probing with a load step instead does not
help---under the same conditions it gives $N=(0.13,0.00,0.00)$, less
informative still. And over $497$ runs on four process-typical plants,
one-signed records never produced a count anywhere near the limits of
\cite{SPIN}: their maximum was $0.32$ of the limit while true $M_s$
ranged from $1.2$ to $6.0$. The reading is one-sided. It can say the
peak is at least so high; it cannot say the peak is low.

\section{Measurements}
\label{sec:measure}

Every number here tests a prediction made above. Nothing is fitted.

\paragraph{The identity, with no adjustable constant.}
Equation~\eqref{eq:law} says the dimensionless ratio
\begin{equation}
R=\frac{M_s\,c\sqrt{n_m^2+n^2}}{N}
\label{eq:ratio}
\end{equation}
equals one for a marginal loop, and approaches one as $M_s$ grows.
Sweeping lag chains of order one to four with dead time, integral times
$T_i/T\in\{0.5,1,2\}$ and the gain swept:

\begin{center}
\begin{tabular}{lccccc}
\toprule
$M_s$ & 1--2 & 2--3 & 3--4 & 6--7 & 8--9 \\
\midrule
$R$ (median) & 1.754 & 1.562 & 1.256 & 1.131 & 1.074 \\
interquartile & 1.74--1.78 & 1.32--1.75 & 1.24--1.27 & 1.12--1.13 & 1.07--1.08 \\
\bottomrule
\end{tabular}
\end{center}

\noindent
$R$ falls monotonically towards one, reaching $1.074$ at $M_s\approx8.5$
and still decreasing. This is the behaviour \eqref{eq:law} predicts: the
identity is exact for a marginal loop, and the excess at moderate $M_s$
is the baseline correction noted above. The interquartile ranges are a
percent or two wide, so the agreement is not an average over scatter.

\paragraph{Only the slopes matter.}
Equation~\eqref{eq:law} claims the plant enters through $n$ and $n_m$
alone. Breaking the same sweep out by plant order tests it directly:

\begin{center}
\begin{tabular}{lccc}
\toprule
$M_s$ range & order 1 & order 2 & order 3 \\
\midrule
2--3 & 1.549 & 1.502 & 1.384 \\
3--5 & 1.185 & 1.189 & 1.284 \\
5--8 & ---   & 1.106 & 1.154 \\
\bottomrule
\end{tabular}
\end{center}

\noindent
At each level the three orders agree to within a few percent, and the
residual spread is smaller than the drift with $M_s$. Pooling all orders
at $M_s\ge5$ gives $R=1.104$ over a range $1.036$ to $1.176$.

\paragraph{Dead time is invisible, as \eqref{eq:npi2} says.}
Recalibrating the map at each $\tau$ separately and comparing over
$N\in[0.3,1.5]$, the largest shift against the $\tau=0.5$ reference is
$0.041$ at $\tau=0.3$, $0.041$ at $\tau=0.7$ and $0.064$ at $\tau=0.9$.
One calibration serves the range, which is \eqref{eq:npi2} restated as a
measurement.

\paragraph{Readings on fixed tunings.}
Four process-typical plants \cite{SPIN}, PI controllers with
$T_i/L_d\in\{1.5,2,3,4,6\}$ and four gain levels, one setpoint step
each, compared against exact $M_s$. On valid records: $n=24$, median
error $+5\%$, interquartile range $-7\%$ to $+8\%$, and $22$ of $24$
within $\pm25\%$. This is the practical cost of
Remark~\ref{rem:scope}---the working loop is not marginal, and its $T_i$
does not exactly cancel---and it is modest.

\paragraph{The validity condition catches the failures.}
In a matched test over five classical tuning rules, every gross error
---$-51\%$, $-47\%$, $-37\%$, $-31\%$---fell on a record the condition
rejects, while accepted readings stayed within $-10\%$ to $+14\%$.

\section{An application: specifying a tuning rule}
\label{sec:spin}

The \spin rule \cite{SPIN} lowers the gain of the lowest band whose
count exceeds a limit, with limits $\bar N=(0.5,0.75,1.0)$ fixed by
calibration. Reading them through the law gives peaks of
\begin{equation}
(1.71,\;1.92,\;2.29),
\label{eq:spinlimits}
\end{equation}
a graded ladder: two inside the classical design window and the third
just above it. The calibration landed there without being aimed there.

The limits may therefore be \emph{stated} instead of chosen: pick the
admissible peak, invert the law, obtain $\bar N$. The specification is
also finer than usual, since it is per band rather than one global
number.

\section{Conclusion}

A setpoint step already contains the sensitivity function, because
$E(s)=S(s)/s$. Counting how many times the error curls around its
settling point measures the damping of the dominant pole pair, and the
peak follows by a linear law whose only plant-dependent constant is the
phase slope at crossover.

For a PI loop with the lag cancelled, that constant is $\pi/2$ exactly,
whatever the dead time, because a dead time contributes equally to phase
and to phase slope and the crossover condition pins its share. This is
the reason the reading needs no model. Adding derivative action removes
the pinning, and with it the invariance---which is why this paper is
about PI loops.

Two limitations bound the result. The analysis is noise-free, and noise
enters the differenced coordinates directly. And the reading is
one-sided: it is valid when the error changes sign and silent otherwise.
Both are stated conditions rather than hidden ones, and the second is
decided from the same record it qualifies.

\section*{Data availability}

Scripts reproducing every number are included with the \spin reference
implementation
(\url{https://github.com/ottapav/roboPID-simulator}).

\begin{thebibliography}{00}

\bibitem{SPIN} D.~Pachner, P.~Otta, J.~Dost\'al, V.~Havlena, Model-free
PID tuning by step-response inspection, submitted.

\bibitem{AstromHagglund2001} K.J.~\r{A}str\"om, T.~H\"agglund, The future
of PID control, Control Eng.\ Pract.\ 9~(11) (2001) 1163--1175.

\bibitem{AstromHagglund2004} K.J.~\r{A}str\"om, T.~H\"agglund,
Revisiting the Ziegler--Nichols step response method for PID control,
J.~Process Control 14~(6) (2004) 635--650.

\bibitem{Skogestad2003} S.~Skogestad, Simple analytic rules for model
reduction and PID controller tuning, J.~Process Control 13~(4) (2003)
291--309.

\bibitem{Desborough2002} L.~Desborough, R.~Miller, Increasing customer
value of industrial control performance monitoring---Honeywell's
experience, in: Chemical Process Control VI, AIChE Symposium Series
No.~326, Vol.~98, 2002, pp.~169--188.

\bibitem{AstromHagglundBook2006} K.J.~\r{A}str\"om, T.~H\"agglund,
Advanced PID Control, ISA, Research Triangle Park, NC, 2006.

\end{thebibliography}

\end{document}
